\chapter{Introduction}
\label{ch:introduction}

\section{Topology control}
\label{sec:intro-topology-control}

A transmission grid is not a fixed graph. Inside a substation, every connected
element, each transmission line, generator, and load, is attached to one of
several busbars, and an operator can move an element from one busbar to another
or split a substation into two electrically separate nodes. Doing so changes
which paths the power actually flows through, and therefore redistributes
loading across the network. This is topology control, also called busbar
switching or network reconfiguration.
\\
\\
Its appeal is economic. Congestion is classically relieved by redispatching
generation or by curtailing renewable production, both of which cost money at
every intervention, or by building new lines, which costs a great deal once and
takes years. Topology control uses assets that already exist and switching
equipment that is already installed, so its marginal cost is close to
zero~\cite{fisher2008ots}. The energy transition makes this attractive rather
than merely elegant: variable renewable generation makes flows less predictable
and pushes existing corridors closer to their thermal limits, while grid
expansion is slow~\cite{marot2021l2rpn,viebahn2022ai}.
\\
\\
The catch is that topology is hard to exploit. Operators use it, but mostly
through a small repertoire of remedial actions established by experience and
offline study, so a large part of the available flexibility is never
searched~\cite{viebahn2022ai}.

\section{Why it has to be automated}
\label{sec:intro-automation}

Three properties make manual search impractical.
\\
\\
The action space is combinatorial. A substation with $d$ connected elements and
two busbars admits $2^{d}$ assignments, and the number of joint configurations
of a grid is the product over its substations. Even the small IEEE 14-bus
benchmark used throughout this thesis exposes $178$ distinct unitary topology
actions; the 36-substation grid of Chapter~\ref{ch:wcci-transfer} exposes more
than $66\,000$, which is why it has to be reduced offline before an agent can be
trained on it at all (Section~\ref{sec:wcci-action-space-reduction}).
\\
\\
The effect of an action cannot be read off the grid. Power flows are governed by
non-linear AC equations, so the consequence of moving one line onto another
busbar is a global redistribution, not a local change. There is no reliable
shortcut for ranking actions other than simulating them.
\\
\\
Finally, the decision is sequential and time-constrained. An action taken now
changes which contingencies are survivable later, cooldown periods forbid acting
on the same substation twice in quick succession, and an operator has minutes,
not hours. Optimizing a single step greedily is not the same problem as keeping
a grid alive for a day.
\\
\\
These are the conditions under which reinforcement learning is a reasonable
candidate: a sequential decision problem, a simulator that is expensive but
available, no analytic optimum, and a clear success signal --- the grid does not
black out. The Learning to Run a Power Network (L2RPN) competition series
was created by RTE to make exactly this problem
tractable for machine-learning research, and the Grid2Op simulator provides the
environment, the chronics of load and generation, and the operational rules that
go with it~\cite{donnot2020grid2op,marot2021l2rpn}.

\section{Why decentralized agents}
\label{sec:intro-decentralized}

Most published work treats the grid as a single agent controlling everything.
Real grids are not operated that way. Transmission system operators divide the
network into regions, each with its own control centre and its own operators,
who act on their own substations with a detailed view of their own area and a
much coarser view of the rest~\cite{viebahn2022ai}. A decentralized formulation
is therefore closer to how the job is actually done, and to how any automated
assistant would eventually be deployed, as support inside one control room,
not as a single controller over an entire interconnection.
\\
\\
It is also the formulation that scales. A monolithic controller has to represent
the joint action space, which grows multiplicatively with the number of
substations. Partitioning the grid into regions and giving each region its own
policy replaces that product by a sum: each agent chooses among its own local
actions only. This is the classical argument for decentralized multi-agent
reinforcement learning, and it comes with the classical
cost~\cite{tan1993independent}. Because every agent learns while the others are
also changing, each one faces an environment that is non-stationary from its own
point of view, and because each observes only part of the grid, agents can take
individually sensible actions that are jointly harmful.
\\
\\
The usual answer is \emph{centralized training with decentralized execution}: a
critic is allowed to see the global state during training, while each actor
conditions only on what it will have access to at deployment
time~\cite{lowe2017maddpg,yu2022mappo}. This thesis uses that setting throughout.

\section{Reinforcement learning for topology control}
\label{sec:intro-related-work}

Work on RL for topology control falls into four groups, in roughly historical
order. A broader survey of graph-based RL for power systems is given
by~\textcite{hassouna2024survey}.

\paragraph{Single-agent policies.}
The first approaches train one agent to control the whole grid.
\textcite{subramanian2021exploring} showed that a plain deep RL agent acting on a
curated set of topology actions already outperforms a do-nothing policy on
L2RPN benchmarks. The strongest example is
SMAAC~\cite{yoon2021smaac}, which won the L2RPN WCCI 2020 challenge with an
actor-critic method over \emph{afterstates} --- the grid configuration that
results from an action, evaluated before the physics resolves it --- combined
with a hierarchical decomposition of where to act and how. The common obstacle
in this group is the action space: it has to be reduced, curated, or factorized
before learning becomes feasible at all.

\paragraph{Learned policies wrapped in heuristics.}
A second group accepts that a learned policy should not be in charge at every
step, and surrounds it with rules. The dominant pattern is an activation
threshold: the agent is queried only when the grid is stressed, measured by the
maximum line loading $\rho_{\max}$, and does nothing otherwise. PowRL added a
heuristic overload manager and a reduced action set on top of an RL policy and
topped the L2RPN NeurIPS 2020 robustness track~\cite{chauhan2023powrl}.
\textcite{lehna2023comparative} compared advanced rule-based agents against RL
agents directly and found comparable survival, with the RL agent far cheaper at
inference --- a useful reminder that beating a well-engineered heuristic is not
automatic. Heuristics of this kind are now standard practice rather than a
baseline, and this thesis uses one.

\paragraph{Hierarchical and centrally coordinated multi-agent methods.}
A third group decomposes the decision but keeps a coordinator.
\textcite{manczak2023hrl} split control into choosing whether to act, choosing
where, and choosing which local configuration to apply.
\textcite{vandersar2023hmarl} assign a low-level agent to each substation under a
higher-level agent that selects which substation acts.
\textcite{demol2025ccma} make the coordination explicit: regional agents propose
actions and a central agent selects among the proposals. These methods recover
much of the scalability of factorization, but they retain a component that sees
the whole grid and decides for everyone, so they do not answer whether purely
local policies suffice.

\paragraph{Decentralized multi-agent control.}
The remaining question --- whether agents that act on local observations, without
a coordinator and without communication, can operate a grid --- is much less
explored, and it is where this thesis sits. The gap is recognized: the
MARL2Grid-TR benchmark was built specifically because prior work
``focused on single-agent settings, neglecting the often decentralized,
multi-agent nature of grid control''~\cite{marchesini2026marl2grid}.

\paragraph{Graph representations.}
Cutting across all four groups is the question of how the grid is given to the
network. A flat vector ties the policy to one grid size and discards the
topology entirely. Graph neural networks are the natural
alternative~\cite{kipf2017gcn,velickovic2018gat,xu2018how}, and their parameters
are independent of the number of substations, which is what makes transfer
between grids conceivable at all. The choice of representation is not neutral,
however: \textcite{dejong2025gnntopo} show that the standard homogeneous graph
suffers from \emph{busbar information asymmetry} --- it encodes the connections
that currently exist but not the ones an action could create --- and that a
heterogeneous representation resolves it. That observation is close to the
subject of Chapters~\ref{ch:methods} and~\ref{ch:experiments}.

\section{Starting point: MARL2Grid and MAPPO}
\label{sec:intro-marl2grid}

This thesis builds on MARL2Grid, the multi-agent benchmark and codebase of
\textcite{marchesini2026marl2grid}, itself the multi-agent successor of
RL2Grid~\cite{marchesini2025rl2grid}. It wraps Grid2Op in a multi-agent
interface and supplies the decomposition, the observation and action spaces, the
reward, and the training loop that the rest of this work modifies.

\paragraph{Decomposition.}
The grid is partitioned into disjoint regions, one per agent, fixed in advance.
On the IEEE 14-bus environment (\texttt{bus14}, 14 substations and 20 lines)
three agents own substations $\{0,1,2,4\}$, $\{3,6,7,8\}$ and
$\{5,9,\ldots,13\}$, giving them $52$, $50$ and $76$ non-idle actions
respectively --- the $178$ actions of the full grid, partitioned rather than
multiplied. Each agent observes only quantities belonging to its own region and
chooses among the topology actions available at its own substations; action $0$
is always do-nothing. Agents act simultaneously and are given no information
about what the others are doing.

\paragraph{Reward and heuristic.}
The training reward is dominated by a flat $+1$ per surviving step, plus an
overload term, so the objective is essentially to keep the episode alive; the
reported metric is episodic survival, the fraction of evaluation episodes
completed without a blackout. Following the practice described above, an idle
heuristic fast-forwards the simulation while the grid is safe: agents are
queried only once $\rho_{\max} \geq 0.90$.

\paragraph{MAPPO.}
The learning algorithm is multi-agent PPO~\cite{yu2022mappo}, the natural
multi-agent extension of PPO~\cite{schulman2017ppo} under centralized training
with decentralized execution. Each agent $i$ has its own stochastic policy
$\pi_{\theta_i}(a_i \mid o_i)$ over its discrete local action space, and is
updated with the clipped surrogate objective
\[
  L^{\text{CLIP}}(\theta_i)
  = \hat{\mathbb{E}}_t\!\left[
      \min\!\Big(
        r_t(\theta_i)\,\hat{A}_t,\;
        \operatorname{clip}\big(r_t(\theta_i),\,1-\epsilon,\,1+\epsilon\big)\,
        \hat{A}_t
      \Big)
    \right],
  \qquad
  r_t(\theta_i)
  = \frac{\pi_{\theta_i}(a_{i,t} \mid o_{i,t})}
         {\pi_{\theta_i^{\text{old}}}(a_{i,t} \mid o_{i,t})},
\]
where $\hat{A}_t$ is a generalized advantage estimate. The advantages come from
a single \emph{centralized} value function trained on the global state, which is
available in simulation but never given to the actors. Only the actors are kept
at execution time, so the deployed system is decentralized even though training
was not.

\paragraph{What this thesis starts from.}
The baseline is that system with flat multilayer-perceptron actors and critic,
trained on \texttt{bus14}. It is unstable at low survival out of the box, and it
is tied to one grid: an MLP's input width is a function of the number of
substations, so nothing it learns can be reused elsewhere. The chapters that
follow address these limitations in the order they arose: first stabilize
training, then make interventions more operator-like, and finally replace the
fixed-width representation.

\section{Outline}
\label{sec:intro-outline}

Chapter~\ref{ch:baseline} tunes the flat baseline until it trains stably.
Chapter~\ref{ch:sparse-control} then changes the objective and asks whether the
same agents can keep their survival while intervening rarely; its preliminary
WCCI pilot exposes the limitations of retaining a fixed-width MLP actor.
Chapter~\ref{ch:methods} consequently defines size-independent graph inputs:
feature scaling, three grid representations, substation-level augmentations,
the readout, and the candidate-action head. Chapter~\ref{ch:experiments} first
checks that a shared graph encoder remains competitive, then screens those
design choices. Chapter~\ref{ch:wcci-transfer} moves to the 36-substation WCCI
grid and tests what actually transfers.
